% Clase 2 — Búsqueda en Ambientes Complejos
% Lectura: AIMA Cap. 4 (Hill Climbing + Algoritmos Genéticos)
\section{Búsqueda en Ambientes Complejos}

\subsection{Búsqueda Local}

A diferencia de la búsqueda sistemática, la búsqueda local opera sobre un \textbf{estado actual} sin mantener el camino recorrido. Útil cuando el objetivo es el estado mismo (no el camino).

\subsection{Hill Climbing (Ascenso de Colina)}

\begin{algorithm}[H]
\caption{Hill-Climbing}
\begin{algorithmic}[1]
\State $\text{current} \leftarrow \text{estado inicial}$
\Loop
    \State $\text{neighbor} \leftarrow$ mejor sucesor de $\text{current}$
    \If{$f(\text{neighbor}) \leq f(\text{current})$}
        \Return $\text{current}$
    \EndIf
    \State $\text{current} \leftarrow \text{neighbor}$
\EndLoop
\end{algorithmic}
\end{algorithm}

\textbf{Problemas}:
\begin{itemize}
    \item Máximos locales (no globales)
    \item Mesetas (\textit{plateaus}) y crestas
\end{itemize}

\textbf{Variantes}: Hill climbing estocástico, reinicio aleatorio, \textit{simulated annealing}.

\subsection{Simulated Annealing}

Permite movimientos malos con probabilidad $e^{\Delta E / T}$ donde $T$ decrece con el tiempo. Garantiza convergencia al óptimo global si $T$ baja suficientemente despacio.

\subsection{Algoritmos Genéticos}

Inspirados en la evolución biológica. Operan sobre una \textbf{población} de individuos (estados codificados como cadenas).

Operadores:
\begin{enumerate}
    \item \textbf{Selección}: individuos con mayor aptitud (\textit{fitness}) tienen más probabilidad de reproducirse.
    \item \textbf{Cruzamiento} (\textit{crossover}): combina dos individuos para producir descendencia.
    \item \textbf{Mutación}: cambia aleatoriamente genes con baja probabilidad.
\end{enumerate}

% Agrega aquí ejemplo del problema de las N-reinas con AG
